% Deep Learning — Topic proposal (Milestone 0)
\documentclass[10pt,a4paper]{article}
\usepackage[utf8]{inputenc}
\usepackage[T1]{fontenc}
\usepackage{lmodern}
\usepackage[margin=2cm]{geometry}
\usepackage{microtype}
\usepackage{xcolor}
\usepackage[numbers,sort&compress]{natbib}
\usepackage{hyperref}
\hypersetup{
  colorlinks=true,
  linkcolor=blue!50!black,
  citecolor=blue!50!black,
  urlcolor=blue!65!black
}
\usepackage{enumitem}
\usepackage{amsmath,amssymb}
\setlist{nosep,leftmargin=*}
\setlength{\parskip}{0.32em}
\setlength{\parindent}{0pt}

\title{\textbf{FB15k-237 Link Prediction:\\Matched KGE, Custom Hard Negatives, and Relation-Level Analysis}}
\author{%
  \begin{tabular}{rl}
    Thalia Delhaise  & student no.\ \textit{66937} \\[0.2em]
    Lenny Briclet    & student no.\ \textit{66926} \\[0.2em]
    André Fonseca    & student no.\ \textit{58225} \\[0.2em]
    Rafael Gufler    & student no.\ \textit{66081}
  \end{tabular}\\[0.45em]
  \small Deep Learning \& Knowledge Graphs --- unified project (all four members).%
}
\date{}

\begin{document}
\maketitle
\thispagestyle{empty}
\vspace{-1.1em}

\textbf{Problem and motivation.}
A \emph{knowledge graph} collects facts as triples $(h,r,t)$. \emph{Link prediction} scores queries $(h,r,?)$ and $(?,r,t)$ using a learned $f_\theta(h,r,t)\in\mathbb{R}$~\citep{bordes2013translating}. \textbf{Research question:} under \textbf{matched} training conditions on FB15k-237, where do \textbf{\textsc{RotatE} vs.\ \textsc{TransE}} differ in \textbf{filtered} ranking---only in aggregate, or mainly on \textbf{specific relation / degree regimes}? Second, does \textbf{custom hard-negative training} (our PyTorch code) change those gaps compared to \textbf{default} library corruption? \textbf{What we aim to learn:} \textbf{quantified} trade-offs (frequent vs.\ rare relations; default vs.\ hard negatives), not a new SOTA. A shallow \textbf{relational GNN} (R-GCN / CompGCN) is \textbf{optional} for the final report if time remains; the core is KGE $+$ custom negatives $+$ \textbf{slice} and \textbf{error} analysis.

\textbf{Pipeline (build the KG, then predict).}
The project matches both courses in order: \textbf{(1)~Dataset $\rightarrow$ graph.} We use the public \textbf{FB15k-237} benchmark~\citep{sun2019rotate} - released as \textbf{real triple files} (train/validation/test splits). That release is our dataset: thousands of facts about real-world entities. \textbf{Building} the KG means ingesting those triples, mapping strings to entity/relation IDs, forming the edge set $\mathcal{T}$ and the split, and loading the result into PyKEEN~\citep{ali2021pykeen}; we may also export a small RDF or edge-list slice for figures or the KG report. \textbf{(2)~Prediction.} We then train models only on $\mathcal{T}_{\mathrm{train}}$ and evaluate on held-out test triples (next paragraph). We do not invent a private ontology from scratch, but we still \textbf{materialize} the full multi-relational graph the methods run on.

\textbf{What the graph looks like.}
Nodes are entities; each fact is a \textbf{directed, labeled} link $h \xrightarrow{r} t$. Many relation types exist, and the same pair $(h,t)$ can appear under different $r$. Equivalently, $\mathcal{G}=(\mathcal{E},\mathcal{R},\mathcal{T})$ with $\mathcal{T}\subseteq\mathcal{E}\times\mathcal{R}\times\mathcal{E}$. We keep the official benchmark partition unchanged (no extra subsampling).

\textbf{Task and evaluation protocol.}
Training uses only $\mathcal{T}_{\mathrm{train}}$: we minimize a \textbf{margin/ranking loss} and draw \textbf{negative} triples by corrupting head or tail (standard pipeline~\citep{bordes2013translating,sun2019rotate}). At test time, each $(h,r,t)\in\mathcal{T}_{\mathrm{test}}$ was unseen during training; we score all candidate completions and record the rank of the true entity. \textbf{Filtered} metrics~\citep{bordes2013translating}: when ranking the true tail for $(h,r,?)$, any other tail $t'$ with $(h,r,t')\in\mathcal{T}_{\mathrm{train}}\cup\mathcal{T}_{\mathrm{val}}\cup\mathcal{T}_{\mathrm{test}}$ is treated as ambiguous (not penalized as an error); the same for head prediction. We average results over head and tail prediction.

\textbf{Scale and Milestone~1 characterization.}
FB15k-237 counts: \textbf{14{,}541} entities, \textbf{237} relations, \textbf{272{,}115} / \textbf{17{,}535} / \textbf{20{,}466} train/val/test triples. By Milestone~1 we add \textbf{dataset statistics} beyond counts: relation-frequency and (head/tail) \textbf{degree distributions}, plus short examples of \textbf{one-to-one / one-to-many / many-to-one} style patterns to motivate later error analysis. Optional second graph (e.g.\ OGB); FB15k-237 stays primary.

\textbf{Method (planned).}
\textbf{Core comparison (matched where possible):} \textsc{TransE}~\citep{bordes2013translating} vs.\ \textsc{RotatE}~\citep{sun2019rotate} with the \textbf{same} splits, training budget, embedding size $d$, and filtered evaluator; PyKEEN~\citep{ali2021pykeen} for reproducible \textbf{data loading} and \textbf{default} KGE runs. \textbf{Controlled ablation axis:} negative-sampling / corruption strategy (default vs.\ our extension in~(1)); if time, a second axis (e.g.\ $d$ or training epochs) for efficiency--accuracy curves.

\textbf{Metrics, relation-wise analysis, and error study.}
\textbf{Aggregate:} filtered MRR and Hits@$1$, Hits@$3$, Hits@$10$ (head, tail, then averaged). \textbf{Relation-wise:} same metrics \textbf{sliced} by relation frequency or degree bins (addresses ``only a single number'' risk). \textbf{Qualitative error analysis:} inspect successes/failures and relate them to relation ambiguity, tail/head frequency, or pattern type. \textbf{Conceptual note:} we keep \textbf{filtered} ranking as primary and briefly contrast why raw ranks would be optimistic~\citep{bordes2013translating}. Training/validation curves for all models.

\textbf{Our contribution vs.\ library defaults.}
PyKEEN gives a fair \textbf{baseline pipeline} (data + default KGE). Our \textbf{own code} (PyTorch) implements the \textbf{custom hard-negative} module, optional \textbf{bootstrap} for CIs, and \textbf{slice} evaluation scripts. The project is not ``only two API calls.''

\textbf{Additional experiments (DL depth)---priority order.}
\textbf{(1)~Custom negative sampling.} Learned scorer $g_\phi$ and/or periodic hard-negative mining; \textbf{ablate} against uniform/Bernoulli default. \textbf{(2)~Slice + error analysis.} Relation/degree slices; short qualitative review of failures; optional bootstrap CIs. \textbf{(3)~Optional:} shallow \textsc{R-GCN} or \textsc{CompGCN} on the \textbf{same} splits if schedule allows. \textbf{Main comparison:} default KGE vs.\ KGE with custom negatives, plus relation-wise tables.

\vspace{0.15em}
\begin{footnotesize}
\bibliographystyle{abbrvnat}
\bibliography{references}
\end{footnotesize}

\end{document}
